\documentclass[aspectratio=169, 14pt]{beamer}
\usepackage{stampfly_slides}

\lesson{10}{Python SDK プログラム飛行}{Python SDK Flight}

\begin{document}

% ------------------------------------------------------------------
\begin{frame}
  \titlepage
\end{frame}

% ------------------------------------------------------------------
\begin{frame}{今日のゴール / Today's Goal}
  \begin{block}{目標}
    Python SDK でプログラム飛行を実現する
  \end{block}

  \vspace{1em}

  \begin{itemize}
    \item WiFi 経由で PC から StampFly を制御
    \item SDK API で離陸・移動・着陸をプログラム化
    \item RC 制御値の送信とテレメトリの取得
  \end{itemize}
\end{frame}

% ------------------------------------------------------------------
\begin{frame}{SDK アーキテクチャ / SDK Architecture}
  \centering
  \resizebox{0.85\linewidth}{!}{\documentclass[border=10pt]{standalone}
\usepackage{tikz}
\usetikzlibrary{arrows.meta, positioning, shapes.geometric, calc, fit}

\begin{document}

\definecolor{blockcolor}{RGB}{0, 120, 200}
\definecolor{arrowcolor}{RGB}{80, 80, 80}
\definecolor{wificolor}{RGB}{40, 160, 40}
\definecolor{wscolor}{RGB}{220, 50, 30}

\begin{tikzpicture}[
    block/.style={
      rectangle, draw=blockcolor, line width=1.5pt,
      fill=blockcolor!8, rounded corners=3pt,
      minimum width=32mm, minimum height=18mm,
      font=\bfseries, text=blockcolor, align=center
    },
    wifi/.style={
      rectangle, draw=wificolor, line width=1.5pt,
      fill=wificolor!8, rounded corners=3pt,
      minimum width=28mm, minimum height=18mm,
      font=\bfseries, text=wificolor, align=center
    },
    conn/.style={-{Stealth[length=3mm]}, line width=1.4pt},
    biconn/.style={{Stealth[length=3mm]}-{Stealth[length=3mm]}, line width=1.4pt},
    lbl/.style={font=\small\bfseries, fill=white, inner sep=2pt},
  ]

  % Nodes
  \node[block] (pc) at (0, 0) {PC\\(Python SDK)};
  \node[wifi]  (wifi) at (6, 0) {WiFi AP\\192.168.4.1};
  \node[block] (sf) at (12, 0) {StampFly\\ESP32-S3};

  % TCP CLI connection (top): center of top edges
  \draw[biconn, blockcolor]
    (pc.north) -- ++(0, 1.2)
    -- node[lbl, above] {TCP CLI (port 23)} ($(sf.north)+(0, 1.2)$)
    -- (sf.north);

  % WebSocket connection (bottom): center of bottom edges
  \draw[biconn, wscolor]
    (pc.south) -- ++(0, -1.2)
    -- node[lbl, below, text=wscolor] {WebSocket (port 81)} ($(sf.south)+(0, -1.2)$)
    -- (sf.south);

  % Labels for connections
  \node[font=\scriptsize, blockcolor, right=2mm] at (pc.north) {commands};
  \node[font=\scriptsize, wscolor, right=2mm] at (pc.south) {telemetry 400\,Hz};

  % Internal labels
  \node[font=\scriptsize, arrowcolor, right=2mm] at (sf.south) {sensors / motors};

\end{tikzpicture}
\end{document}
}

  \vspace{0.5em}

  \begin{columns}[c]
    \begin{column}{0.48\textwidth}
      \begin{block}{TCP CLI (port 23)}
        コマンド送受信（takeoff, land...）
      \end{block}
    \end{column}
    \begin{column}{0.48\textwidth}
      \begin{block}{WebSocket (port 81)}
        テレメトリ受信（400\,Hz RT）
      \end{block}
    \end{column}
  \end{columns}
\end{frame}

% ------------------------------------------------------------------
\begin{frame}{Python SDK API}
  \small
  \begin{tabular}{lll}
    \hline
    \textbf{関数} & \textbf{説明} & \textbf{備考} \\
    \hline
    \api{StampFly(host)} & インスタンス生成 & デフォルト 192.168.4.1 \\
    \api{connect()} & WiFi 接続 & TCP CLI + WS \\
    \api{takeoff()} & 離陸 & ブロッキング \\
    \api{land()} & 着陸 & ブロッキング \\
    \api{move\_forward(cm)} & 前進 & 20--200\,cm \\
    \api{rotate\_clockwise(deg)} & 時計回り回転 & 1--360° \\
    \api{send\_rc\_control(lr,fb,ud,yaw)} & RC 制御値送信 & $-100$\textasciitilde$+100$ \\
    \api{get\_telemetry()} & テレメトリ取得 & 400\,Hz dict \\
    \api{end()} & 切断 & リソース解放 \\
    \hline
  \end{tabular}
\end{frame}

% ------------------------------------------------------------------
\begin{frame}[fragile]{実習1: 基本プログラム飛行 / Basic Flight}
  \begin{lstlisting}[style=sfpython, title={flight\_basic.py}]
from stampfly import StampFly

drone = StampFly()
drone.connect()

drone.takeoff()
drone.move_forward(50)    # 50 cm forward
drone.rotate_clockwise(90)
drone.move_forward(50)
drone.land()

drone.end()
  \end{lstlisting}
\end{frame}

% ------------------------------------------------------------------
\begin{frame}[fragile]{実習2: RC制御 + テレメトリ / RC + Telemetry}
  \begin{lstlisting}[style=sfpython, title={flight\_rc.py}]
import time
from stampfly import StampFly

drone = StampFly()
drone.connect()
drone.takeoff()

# Hover + read telemetry for 5 seconds
for _ in range(50):
    drone.send_rc_control(0, 0, 0, 0)
    t = drone.get_telemetry()
    print(f"alt={t.get('alt',0):.1f}")
    time.sleep(0.1)

drone.land()
drone.end()
  \end{lstlisting}
\end{frame}

% ------------------------------------------------------------------
\begin{frame}{チェックポイント / Checkpoint}
  \begin{alertblock}{確認事項}
    \begin{itemize}
      \item[$\square$] WiFi AP に接続し \api{connect()} が成功する
      \item[$\square$] \api{takeoff()} → \api{move\_forward()} → \api{land()} が動作する
      \item[$\square$] \api{get\_telemetry()} で高度データを取得できる
    \end{itemize}
  \end{alertblock}

  \vspace{1em}

  \begin{block}{次のレッスン}
    Lesson 11: ホバリングタイム競技会 ルール説明
  \end{block}
\end{frame}

\end{document}
